\clearpage
\begin{landscape}
\definecolor{donegray}{RGB}{120,133,137}
\definecolor{reviewblue}{RGB}{73,132,174}
\definecolor{sensorsorange}{RGB}{205,139,56}
\definecolor{archteal}{RGB}{41,145,132}
\definecolor{buildgreen}{RGB}{89,151,100}
\definecolor{evalgold}{RGB}{200,166,63}
\begin{figure}[H]
\centering
\small
\textbf{Cronograma de execução e marcos do projeto}\par
\vspace{1mm}
\footnotesize Outubro de 2026 a março de 2028\par
\vspace{2mm}
\begin{tikzpicture}[x=0.68cm,y=0.84cm]


% Months: Oct--Dec/2026, Jan--Dec/2027, Jan--Mar/2028.
\foreach \i/\mes in {0/Out,1/Nov,2/Dez,3/Jan,4/Fev,5/Mar,6/Abr,7/Mai,8/Jun,9/Jul,10/Ago,11/Set,12/Out,13/Nov,14/Dez,15/Jan,16/Fev,17/Mar}{%
    \pgfmathsetmacro{\monthx}{\i+0.5}%
    \node[rotate=90,font=\tiny,anchor=center] at (\monthx,1.35) {\mes};%
}
\node[font=\scriptsize\bfseries] at (1.5,1.85) {2026};
\node[font=\scriptsize\bfseries] at (9,1.85) {2027};
\node[font=\scriptsize\bfseries] at (16,1.85) {2028};
\node[anchor=east,font=\scriptsize\bfseries] at (-0.3,1.35) {ATIVIDADE / MARCO};

\foreach \i in {0,...,18}{%
    \draw[gray!25] (\i,1.1) -- (\i,-11.35);%
}
\foreach \i in {0,...,29}{%
    \pgfmathsetmacro{\gridy}{0.95-0.42*\i}%
    \draw[gray!18] (0,\gridy) -- (18,\gridy);%
}
\draw[gray!65,thick] (3,1.1) -- (3,-11.35);
\draw[gray!65,thick] (15,1.1) -- (15,-11.35);

\newcommand{\barrow}[5]{%
    \pgfmathsetmacro{\rowy}{0.76-0.42*#1}%
    \node[anchor=east,align=right,text width=10.15cm,font=\tiny] at (-0.3,\rowy) {#4};%
    \fill[#5] (#2,{\rowy-0.12}) rectangle (#3,{\rowy+0.12});%
}

\barrow{0}{0}{1}{Busca bibliográfica e triagem inicial: 135 registros; 15 publicações para análise preliminar.}{donegray}
\barrow{1}{1}{4}{Leitura integral, confirmação da elegibilidade e atualização da seleção.}{reviewblue}
\barrow{2}{2}{6}{Extração dos estudos, análise crítica e síntese do estado da arte.}{reviewblue}
\barrow{3}{1}{6}{Revisão dos capítulos de introdução e fundamentação teórica.}{reviewblue}
\barrow{4}{1}{4}{Definição do local, permissões e plano de instalação dos sensores adquiridos.}{sensorsorange}
\barrow{5}{3}{6}{Instalação, configuração, calibração e testes iniciais dos sensores.}{sensorsorange}
\barrow{6}{3}{6}{Preparação da ingestão, armazenamento temporal e controle de qualidade.}{archteal}
\barrow{7}{5}{12}{Coleta contínua, inspeção, manutenção e calibração dos sensores.}{sensorsorange}
\barrow{8}{6}{11}{Limpeza, padronização e integração das séries com dados contextuais.}{archteal}
\barrow{9}{3}{6}{Comparação de arquiteturas e definição de requisitos e casos de uso.}{archteal}
\barrow{10}{4}{7}{Especificação dos componentes, agentes, fontes e mecanismos de rastreabilidade.}{archteal}
\barrow{11}{5}{8}{Seleção do detector e avaliação da linha de base analítica.}{archteal}
\barrow{12}{6}{9}{Implementação dos serviços de consulta a dados e análise de eventos.}{buildgreen}
\barrow{13}{7}{10}{Curadoria da base documental e implementação da recuperação RAG.}{buildgreen}
\barrow{14}{7}{10}{Implementação do orquestrador, integração com LLM e interface de consulta.}{buildgreen}
\barrow{15}{9}{11}{Integração ponta a ponta, registros, tratamento de falhas e controle de acesso.}{buildgreen}
\barrow{16}{7}{9}{Definição do protocolo experimental, cenários e métricas de avaliação.}{evalgold}
\barrow{17}{9}{11}{Testes funcionais, de integração, confiabilidade, latência e custo.}{evalgold}
\barrow{18}{10}{12}{Experimentos com dados coletados e cenários complementares controlados.}{evalgold}
\barrow{19}{11}{13}{Análise dos resultados, limitações e validade dos experimentos.}{evalgold}
\barrow{20}{7}{12}{Redação do texto e preparação dos resultados parciais para qualificação.}{reviewblue}
\barrow{21}{12}{13}{Exame de qualificação, com data-limite de 27/10/2027.}{evalgold}
\barrow{22}{13}{14}{Incorporação das recomendações da banca de qualificação.}{reviewblue}
\barrow{23}{12}{15}{Redação e submissão de artigo científico com os resultados da pesquisa.}{reviewblue}
\barrow{24}{12}{15}{Redação dos capítulos de metodologia, implementação, resultados e discussão.}{reviewblue}
\barrow{25}{14}{16}{Revisão integral, normalização e encaminhamento da versão à banca.}{reviewblue}
\barrow{26}{15}{17}{Preparação da apresentação e dos materiais para defesa.}{reviewblue}
\barrow{27}{16}{17}{Defesa da dissertação, prevista para fevereiro de 2028.}{evalgold}
\barrow{28}{16}{18}{Correções da banca e depósito da versão final até março de 2028.}{reviewblue}

\end{tikzpicture}

\vspace{1mm}
\scriptsize
\colorbox{donegray}{\rule{0.35cm}{0.1cm}} Concluído/documentado\quad
\colorbox{reviewblue}{\rule{0.35cm}{0.1cm}} Revisão e redação\quad
\colorbox{sensorsorange}{\rule{0.35cm}{0.1cm}} Sensores e dados\quad
\colorbox{archteal}{\rule{0.35cm}{0.1cm}} Arquitetura\quad
\colorbox{buildgreen}{\rule{0.35cm}{0.1cm}} Desenvolvimento\quad
\colorbox{evalgold}{\rule{0.35cm}{0.1cm}} Avaliação e marcos

\caption{Cronograma mensal e marcos previstos para o desenvolvimento da pesquisa.}
\label{fig:cronograma-pesquisa}
\fonte{Elaboração própria. A busca e a triagem inicial foram realizadas em outubro de 2026. A qualificação deverá ocorrer até 27 de outubro de 2027; a defesa deverá ocorrer até fevereiro de 2028 e o depósito da versão final até março de 2028, respeitados os prazos do programa.}
\end{figure}
\end{landscape}
\clearpage
